\section*{Limitations}

Despite the effectiveness and practicality demonstrated by QEVA, our proposed metric inherits several limitations inherent to its underlying models and design.

\paragraph{Hallucination.} 
As QEVA employs a Large Multimodal Model (LMM) for question generation and answering, it may occasionally produce hallucinated content in the generated questions or answers not actually present in the video. Although our filtering process significantly mitigates this issue, the possibility of subtle hallucinations remains, potentially affecting evaluation reliability in edge cases.

\paragraph{API Cost and Processing Speed.}
QEVA relies heavily on inference from Large Language Models (LLMs) and Video-LMMs. Such models typically require significant computational resources and incur relatively high API costs, particularly when evaluating large-scale datasets or numerous summaries. This dependence may limit QEVA's practical applicability in resource-constrained environments or real-time scenarios.

\paragraph{Necessity of Post-processing.}
QEVA occasionally produces outputs that deviate from the predefined format or scoring criteria. Although infrequent, these cases necessitate additional post-processing to ensure compliance with the intended evaluation guidelines, slightly complicating the evaluation pipeline.

\paragraph{Preference for LLM-based Outputs.}
Recent evaluations using LLMs have identified a subtle preference bias towards outputs generated by LLMs themselves. QEVA may similarly exhibit a slight bias favoring summaries produced by certain Video-LMMs, potentially influencing the fairness and objectivity of the evaluation. This phenomenon warrants further investigation to quantify and mitigate such biases in future research.

\section*{Ethics Statement}
Ethics Statement
Our research involved a user study with human participants to collect judgments for evaluating our proposed metric. The study was conducted in accordance with established ethical guidelines. We recruited 20 annotators who participated on a voluntary basis. While our institution's review board determined that formal IRB approval was not required for this type of user study (involving subjective evaluation of anonymized system outputs), we took several measures to ensure the protection of participants.

Prior to the study, all participants were provided with a clear description of the research objectives and the annotation task, and we obtained informed consent from each individual. To acknowledge their valuable contribution, participants were fairly compensated for their time and effort. All data collected during the study was fully anonymized to protect the privacy and confidentiality of the participants, and no personally identifiable information (PII) was collected or stored.

\section*{Acknowledgement}
This work was supported by IITP grant funded by the Korea government(MSIT) (No. RS-2020-II200004, Development of Previsional Intelligence based on Long-term Visual Memory Network) and supported by the Institute of Information \& Communications Technology Planning \& Evaluation (IITP) grand funded by the Korea government (MSIT) [RS-2021-II211341, Artificial Intelligence Graduate School Program(Chung-Ang University)].